% ECONOMICS_PROBLEM: {"id": "J41-611", "title": "Temporary contracts: stepping stones or traps", "jel_code": "J41", "source_id": "OP-0359"}
% !TeX program = xelatex
\documentclass[11pt,a4paper]{article}
\usepackage[margin=23mm]{geometry}
\usepackage{fontspec}
\setmainfont{Latin Modern Roman}
\usepackage{amsmath,amssymb,mathtools}
\usepackage{xcolor,enumitem,booktabs,longtable,array,xurl,hyperref}
\definecolor{muted}{HTML}{53636C}
\hypersetup{colorlinks=true,urlcolor=blue,linkcolor=blue}
\setlength{\parindent}{0pt}
\setlength{\parskip}{5pt}
\newcommand{\R}{\mathbb R}
\newcommand{\E}{\mathbb E}
\newcommand{\N}{\mathbb N}
\newcommand{\ind}{\mathbf 1}
\newcommand{\Var}{\operatorname{Var}}
\newcommand{\Cov}{\operatorname{Cov}}
\newcommand{\supp}{\operatorname{supp}}
\DeclareMathOperator*{\argmin}{arg\,min}
\DeclareMathOperator*{\argmax}{arg\,max}
\newcommand{\entrymeta}[3]{{\sffamily\small\color{muted}\textbf{\texttt{#1}}\quad|\quad #2\par #3\par}}
\newcommand{\Problemref}[1]{\texttt{#1}}
\newcommand{\JELref}[2]{\texttt{#2} (JEL \texttt{#1})}
\begin{document}
\section*{Temporary contracts: stepping stones or traps}
\textbf{Problem ID:} \texttt{J41-611}\quad \textbf{JEL:} \texttt{J41}\par
% BEGIN ECONOMICS STATEMENT
\label{problem:OP-0359}
{\sffamily\small\color{muted}Original subject group: Labor markets and work\par}
\label{definitions:problem:30007641}
\textbf{Audit classification:} Background; proposed extension. Both working papers checked and journal metadata distinguished. They provide findings and a synthesis; the broad causal and reform extensions are editorial.
\subsection*{Definitions and precise research target}
\textbf{Complete editorial problem.} Fix a cohort of jobseekers, a labor market, and an employment horizon \(T\). Let \(D_i=1\) denote acceptance of a precisely specified temporary contract and \(D_i=0\) the declared alternative, such as continued job search rather than immediate permanent employment. Let \(A_{it}(d)\) indicate stable employment and \(Y_{it}(d)\) real earnings at date \(t\). Define \(\tau_A(t,x)=E[A_{it}(1)-A_{it}(0)\mid X_i=x]\) and \(\tau_Y(T,x)=E[\sum_{t=0}^{T}\delta^t\{Y_{it}(1)-Y_{it}(0)\}\mid X_i=x]\), with baseline characteristics \(X_i\) and discount factor \(\delta\in(0,1]\). A stepping stone requires a favorable declared career outcome relative to this comparator, not merely an observed transition to a permanent job. Identify these effects using credible assignment or contract-offer variation, allowing selection and duration dependence. Determine how results vary by contract type, worker disadvantage, demand conditions, and employment-protection rules. Evaluate reforms separately as market-wide policies because firms can substitute among contracts. Earlier studies answer specific settings; the research agenda concerns reliable context-dependent effects and mechanisms, not a conjecture that every temporary job has one common sign.

\textbf{Definitions and admissible scope.} A temporary treatment must specify duration, renewal probability, employer, agency status, pay, and legal protection. The control is a feasible alternative available at the offer date, such as continued search, and must not silently alternate between unemployment and permanent work. Stable employment \(A_{it}\) can mean a permanent contract or paid work in every period of a fixed rolling window; choose one definition. Let \(Z_i\) be randomized offers and \(D_i(Z_i)\) acceptance. An offer effect compares \(A(Z=1)\) with \(A(Z=0)\); a contract-acceptance effect is identified only with additional restrictions. Discounted earnings count zeros and earnings at later employers. Selection conditional on observed permanent employment is not the population effect. Both cited studies already contain answers for particular contracts and contexts. Take a fixed integer \(T\), \(\delta\in(0,1]\), common units for real earnings, and a baseline jobseeker weighting law. The potential outcomes indexed by acceptance \(d\) describe an explicitly feasible contract-versus-search intervention, not a regression comparison between observed acceptors and nonacceptors. Under a randomized offer \(Z\), the directly identified target instead replaces \(d\) by \(Z\); if exclusion, monotonicity, and a nonzero first stage are imposed, the ratio of offer outcome effect to acceptance first stage is a local complier effect. It is not automatically \(\tau_A(t,x)\) or \(\tau_Y(T,x)\) for all jobseekers. Declare whether favorable earnings, stable employment, or a utility criterion defines a stepping stone before evaluating signs.
\subsection*{Variants and assumption changes}
\begin{enumerate}
\item Contract-type comparison (source-informed): compare fixed-term direct hires, agency work, and seasonal jobs against the same search alternative. Estimate outcomes by baseline disadvantage and aggregate demand rather than pooling all temporary contracts.
\item Market reform (editorial): replace individual offers by a reform of dismissal costs or conversion rules. Define equilibrium substitution among contracts, entry, and wages and compare the resulting population welfare with the local offer effect.
\end{enumerate}
\subsection*{References and source audit}
\textbf{Support and access.} Both working papers checked and journal metadata distinguished. They provide findings and a synthesis; the broad causal and reform extensions are editorial. \textbf{Locator:} IZA DP 205 (2000), Introduction pp. 1--2; DP 14367 (2021), Sections 1--2. \textbf{Versions:} Published in The Economic Journal 112(480), F189--F213 (2002), DOI 10.1111/1468-0297.00043; that DOI belongs to the 2002 article. Published under the changed title Are temporary jobs stepping stones or dead ends? A systematic review of the literature, International Journal of Manpower 43(9), 60--74 (2022), DOI 10.1108/IJM-02-2022-0064.  
\begin{enumerate}\raggedright
\item\hypertarget{definitions:source:30007641:1}{} Alison L. Booth; Marco Francesconi; Jeff Frank. 2000. \emph{Temporary Jobs: Stepping Stones or Dead Ends?}. IZA Discussion Paper 205, October 2000, Introduction, pp. 1--2.
\url{https://docs.iza.org/dp205.pdf}.
\item\hypertarget{definitions:source:30007641:2}{} Mattia Filomena; Matteo Picchio. 2021. \emph{Are Temporary Jobs Stepping Stones or Dead Ends? A Meta-Analytical Review of the Literature}. IZA Discussion Paper 14367, May 2021, Sections 1--2.
\url{https://docs.iza.org/dp14367.pdf}.
\item\hypertarget{definitions:source:30007641:3}{} Alison L. Booth; Marco Francesconi; Jeff Frank. 2002. \emph{Temporary Jobs: Stepping Stones or Dead Ends?}. Published abstract; The Economic Journal 112(480), F189--F213.
\url{https://academic.oup.com/ej/article-abstract/112/480/F189/5085411}.
DOI: \url{https://doi.org/10.1111/1468-0297.00043}.
\item\hypertarget{definitions:source:30007641:4}{} Mattia Filomena; Matteo Picchio. 2022. \emph{Are temporary jobs stepping stones or dead ends? A systematic review of the literature}. Published abstract; International Journal of Manpower 43(9), 60--74.
\url{https://www.emerald.com/ijm/article/43/9/60/146599/Are-temporary-jobs-stepping-stones-or-dead-ends-A}.
DOI: \url{https://doi.org/10.1108/IJM-02-2022-0064}.
\end{enumerate}

\subsection*{Common conventions: Source mathematical conventions}

These conventions apply to each entry unless explicitly replaced.

\textbf{Models and identification.} A primitive model \(m\) specifies all probability laws, feasible choices, information, equilibrium and measurement rules used by its target. The admissible class \(\mathcal M\) is fixed independently of outcomes. If \(P_O(m)\) is its observable probability law and \(\tau(m)\) the target, its identified set at observable law \(P\) is
\[
 \mathcal I_\tau(P)=\{\tau(m):m\in\mathcal M,\ P_O(m)=P\}.
\]
Point identification means that this set is a singleton. An empty set signals incompatibility of data and assumptions. Coordinate bounds need not describe the joint set. Bounds are sharp only if every retained target is attainable or arbitrarily approachable by compatible models. Finite bounds require explicit restrictions on unobserved quantities; unrestricted confounding, hidden wealth, extrapolation or arbitrary equilibrium selection can yield uninformative sets.

\textbf{Causal interventions.} A potential outcome \(Y_i(p)\) is the outcome under a complete policy regime \(p\), including timing, financing, information and market spillovers. The target population and units are fixed before treatment. With interference, use the complete assignment vector or a justified exposure mapping; individual no-interference notation alone cannot identify an economy-wide intervention. A difference of mean potential outcomes does not require the cross-world joint distribution. Conditioning on post-treatment migration, survival, enrollment or employment changes the estimand and needs separate justification.

\textbf{Statistical statements.} An estimator, confidence set, forecast or test specifies a sampling law, model class and loss/coverage criterion. Uniform \((1-\alpha)\)-coverage means \(\inf_{m\in\mathcal M}\Pr_m\{\tau(m)\in C_n\}\geq1-\alpha\), or an explicitly asymptotic version. The whole parameter space trivially has coverage, so informativeness requires a stated width, risk or power objective. A forecasting improvement is relative to a named benchmark, information set and evaluation population.

\textbf{Equilibrium and optimization.} An equilibrium comprises feasible plans, optimality, beliefs where needed, budget constraints and all market/resource clearing. The primitives or a stated benchmark define each correspondence \(\mathcal E(p;m)\). Nonemptiness is assumed or proved, never inferred just from compact policy sets. A selected-equilibrium target uses a declared selection rule; a robust target quantifies over all permitted equilibria. Use suprema or infima unless attainment is established. An \(\varepsilon\)-optimal policy is feasible and within \(\varepsilon\) of the finite optimum value. Report an empty feasible set separately.

\textbf{Welfare and accounting.} Normative utility, interpersonal normalization, social weights, numeraire, discounting and the use of public receipts are primitives. Behavioral utility need not coincide with the welfare criterion. Household welfare includes ownership income and rents once; adding profits again to compensating variation generally double-counts them. A monetary equivalent is defined by an explicit transfer and price convention, with monotonicity and range conditions. Average effects, mechanism contributions, transfers and welfare losses are different targets. Infinite sums require convergence and dynamic feasible plans require terminal or no-Ponzi conditions.

\textbf{Source versions.} A working-paper date, revision date and journal publication year are distinguished. A DOI for a journal version is not automatically the DOI of an earlier manuscript. Locators refer to the stated version and printed pages where specified. A metadata-only check does not verify a claim about a full-text conclusion. Every entry's source subsection narrows the provenance claim accordingly.
% END ECONOMICS STATEMENT
\end{document}
